%-*-latex-*-
%-*-latex-*-
\newcommand\COURSE{ciss245}
\newcommand\ASSESSMENT{q2901}
\newcommand\ASSESSMENTTYPE{Quiz}
\newcommand\POINTS{\textwhite{xxx/xxx}}

\input{myquizpreamble}
\input{yliow}
\input{\COURSE}
%-*-latex-*-

%-*-latex-*-

%-*-latex-*-
\renewcommand\TITLE{\ASSESSMENTTYPE \ \ASSESSMENT}

\renewcommand\EMAIL{}
\renewcommand\AUTHOR{eagraham1@cougars.ccis.edu}

\textwidth=6in
\begin{document}
\topmatterthree

%------------------------------------------------------------------------------
\nextq
Compute by hand the value of $f(4)$ if $f(n)$ satisfies the following:
\[
f(n) =
\begin{cases}
1 & \text{if } n = 0 \\
2 & \text{if } n = 1 \\
2 f(n - 2) - 3  & \text{if } n \geq 2
\end{cases}
\]
\ANSWER
\begin{answercode}
-5
\end{answercode}

%------------------------------------------------------------------------------
\nextq
Compute by hand the value of $c$ if $f(4) = 10$
and $f(n)$ satisfies the following:
\[
f(n) =
\begin{cases}
c & \text{if } n = 0 \\
2 & \text{if } n = 1 \\
f(n - 1) + 2 f(n - 2)  + 1  & \text{if } n \geq 2
\end{cases}
\]
Your answer below will be considered a mathematical expression,
not a C\texttt{++} expression.
So when you write $1/3$ below, you mean \lq\lq one third" and not
$0$.
\\
\ANSWER
\begin{answercode}
-5/6
\end{answercode}

%------------------------------------------------------------------------------
\nextq
Write a C\texttt{++} function that implements the following
\[
f(n) =
\begin{cases}
1 & \text{if } n = 0 \\
2 & \text{if } n = 1 \\
3 f(n - 1) + 4 f(n - 2)  & \text{if } n \geq 2
\end{cases}
\]
\ANSWER
\begin{answercode}
int f(int n)
{
  if (n == 0) return 1;
  else if (n == 1) return 2;
  else return 3 * f(n - 1) + 4 * f(n - 2);
}
\end{answercode}

%------------------------------------------------------------------------------
\nextq
Write a C\texttt{++} function that implements the computation of $x^n$ where $x$ is a
\verb!double! and \verb!n! is an integer $\geq 0$.
The following should help with your recursion:
\[
x^{6} = x^{5} \cdot x
\]
(The idea is similar to the recursive factorial computation.)
\\
\ANSWER
\begin{answercode}
double power(double x, int n)
{
  if (n == 0) return 1.0;
  else return power(x, n - 1) * x;
}
\end{answercode}

%------------------------------------------------------------------------------
\nextq
Write a C\texttt{++} function that implements the recursive GCD.
(The one in the notes.)
\\
\ANSWER
\begin{answercode}
int gcd(int a, int b)
{
  if (b == 0) return a;
  else return gcd(b, a % b);
}
\end{answercode}

%------------------------------------------------------------------------------
\nextq
Write a recursive C\texttt{++} function that computes
\[
f(n) = \frac{0}{2^0} + \frac{1}{2^1} + \frac{2}{2^2} + \frac{3}{2^3}
+ \cdots + \frac{n}{2^n}
\]
(You want to think about two things:
What is $f(0)$?
What is the relationship between $f(n+1)$ in terms of $f(n)$?)
\\
\ANSWER
\begin{answercode}
double f(int n)
{
  if (n == 0) return 0.0;
  else return f(n - 1) + n / pow(2.0, n);
}
\end{answercode}

%------------------------------------------------------------------------------
%-*-latex-*-
\newpage
\input{instructions}
\end{document}

